% ===========================================================================
% Section III -- Experimental Setup and Datasets
% Owner: Member 1
% ===========================================================================
\subsection{Bearing Kinematics}
With shaft rotational frequency $f_r=\text{RPM}/60$, rolling-element count
$n_b$, element diameter $d$, pitch diameter $D$ and contact angle $\alpha$,
the ideal characteristic frequencies are
\begin{equation}
\begin{aligned}
  f_{\text{BPFO}} &= \frac{n_b}{2}f_r\left(1-\rho\right), &
  f_{\text{BPFI}} &= \frac{n_b}{2}f_r\left(1+\rho\right),\\
  f_{\text{BSF}} &= \frac{D}{2d}f_r\left(1-\rho^{2}\right), &
  f_{\text{FTF}} &= \frac{f_r}{2}\left(1-\rho\right),
\end{aligned}
  \label{eq:kinematics}
\end{equation}
with the diameter ratio $\rho=(d/D)\cos\alpha$.

The CWRU drive-end bearing is the SKF 6205-2RS: nine rolling elements,
$7.94$~mm element diameter, $39.04$~mm pitch diameter and $0^\circ$ contact
angle. Its calculated frequencies at $1797$~rpm agree with published CWRU
benchmark data~\cite{smith2015}. The MFPT test bearing is described by the
NICE geometry documented for the dataset: eight rolling elements, $0.235$~in
element diameter, $1.245$~in pitch diameter and $0^\circ$ contact angle. The
calculator accepts the geometry explicitly, so the correct bearing model is
used for each dataset.

Table~\ref{tab:kinematics} lists the CWRU frequencies at the four load
settings and the MFPT frequencies at its stated shaft rate of $25$~Hz
($1500$~rpm). These are ideal, geometry-based targets. Actual fault peaks may
be displaced by rolling-element slip and by speed variation, and the detector
of Section~II-D is built to tolerate that displacement.

\begin{table}[!t]
\renewcommand{\arraystretch}{1.15}
\caption{Ideal kinematic fault frequencies calculated from the specified
bearing geometries and shaft speeds}
\label{tab:kinematics}
\centering
\footnotesize
\setlength{\tabcolsep}{4pt}
\begin{tabular}{lccccc}
\toprule
Condition & Speed (rpm) & BPFO & BPFI & BSF & FTF \\
          &             & \multicolumn{4}{c}{(Hz)} \\
\midrule
CWRU 0 HP & 1797 & 107.36 & 162.19 & 70.58 & 11.93 \\
CWRU 1 HP & 1772 & 105.87 & 159.93 & 69.60 & 11.76 \\
CWRU 2 HP & 1750 & 104.56 & 157.94 & 68.74 & 11.62 \\
CWRU 3 HP & 1730 & 103.36 & 156.14 & 67.95 & 11.48 \\
MFPT (25 Hz shaft) & 1500 & 81.12 & 118.88 & 63.86 & 10.14 \\
\bottomrule
\end{tabular}
\end{table}

\subsection{Datasets}
\emph{CWRU:} the Case Western Reserve University Bearing Data
Center~\cite{cwru_datacenter} provides accelerometer recordings for normal
bearings and for seeded inner-race, outer-race and rolling-element defects. A
public NumPy transcription of the archive~\cite{rigas2024} stores the time
series in \texttt{.npz} files with channel keys such as DE, FE and BA. The
filename encodes the shaft-speed family, the defect type, the defect size and,
in many fault files, the drive-end sampling rate through a \texttt{DE12} or
\texttt{DE48} suffix. The loader parses all of it: RPM, fault label, defect
diameter, race position and the rate suffix.

Two cases have to be handled with care. Normal DE files carry no rate suffix,
and the BA rate is never inferred, so both require a verified sampling rate
argument instead of a silent default. The FE acquisition family is $12$~kHz
even when those channels sit alongside \texttt{DE48} data. The load conditions
used here are $1797$~rpm ($0$~HP), $1772$~rpm ($1$~HP), $1750$~rpm ($2$~HP)
and $1730$~rpm ($3$~HP). The broader dataset description sometimes quotes an
overall speed range down to about $1720$~rpm, so the speed of an individual
recording is taken from its metadata where available rather than inferred
from a generic load label~\cite{smith2015}.

\emph{MFPT:} the standard MATLAB recordings of the second benchmark contain a
vibration vector \texttt{gs} together with load, shaft-rate and sample-rate
metadata~\cite{bechhoefer2013}. The varied-load outer-race records cover $25$,
$50$, $100$, $150$, $200$, $250$ and $300$~lb, and the varied-load inner-race
records cover $0$, $50$, $100$, $150$, $200$, $250$ and $300$~lb. The
collection also holds baseline recordings at $270$~lb and $97.656$~kHz for
$6$~s, and fixed-load outer-race records at $270$~lb and the same high rate.
The varied-load fault records are sampled at $48.828$~kHz for $3$~s.

\subsection{Signal Preprocessing}
The first operation subtracts the sample mean,
\begin{equation}
  x_0[n] = x[n] - \frac{1}{N}\sum_{n=0}^{N-1}x[n],
  \label{eq:demean}
\end{equation}
which removes a constant DC offset that would otherwise contribute energy at
zero frequency. Optional peak normalization divides $x_0[n]$ by
$\max|x_0[n]|$, and RMS normalization divides by
$\sqrt{(1/N)\sum_n x_0^2[n]}$. The implementation retains the demeaned signal
as its reference. Normalization is used only when waveform shape is compared;
it should not be used to erase genuine load-dependent amplitude differences
in quantitative physical comparisons.

\subsection{Baseline Spectra and Window Comparison}
The baseline FFT is computed after mean removal and optional windowing. For a
window $w[n]$, the windowed transform is
\begin{equation}
  X_w[k] = \sum_{n=0}^{N-1} x_0[n]\,w[n]\,e^{-j2\pi kn/N}.
  \label{eq:dft}
\end{equation}
A coherent-gain-corrected one-sided amplitude spectrum is obtained by dividing
the magnitude by $\sum_n w[n]$ and doubling the positive-frequency bins other
than DC and, for even $N$, the Nyquist bin. The STFT repeats the windowed
transform over successive overlapping frames, so that the frequency content
can be followed in time.

A finite record may not contain an integer number of cycles of a sinusoidal
component. Multiplication by the observation window then spreads energy into
adjacent DFT bins, which is spectral leakage. The rectangular window has a
narrow main lobe but comparatively high sidelobes; the Hann and Blackman
windows reduce the sidelobes at the cost of a wider main lobe. The
implementation compares all three on the same input and with the same
sampling settings, both for a controlled off-bin sinusoid and for the STFT of
a record.

The controlled test uses $f_s=12$~kHz, $N=2048$ and an off-bin tone at
$1467.01$~Hz. The metric in Table~\ref{tab:leakage} is the fraction of FFT
energy that falls outside an approximate, window-specific main-lobe region.
It is an educational comparison, not a universal standard.

\begin{table}[!t]
\caption{Approximate leakage-energy fraction for a controlled off-bin
sinusoid}
\label{tab:leakage}
\centering
\footnotesize
\begin{tabular}{lr}
\toprule
Window & Leakage-energy fraction \\
\midrule
Rectangular & 0.06613550 \\
Hann        & 0.00017592 \\
Blackman    & 0.00000295 \\
\bottomrule
\end{tabular}
\end{table}

\subsection{Pipeline Parameters}
Unless stated otherwise the experiments use $f_s=12$~kHz, a $4^{\text{th}}$
order Butterworth bandpass realised in cascade SOS form with forward--backward
filtering, an FFT length $N=4096$ (giving $\Delta f = 2.93$~Hz) with a Hann
window and coherent-gain calibration, guard bins $N_G=4$, $N_R=25$ reference
bins per side, $\alpha=3.5$ for CA-CFAR and $\gamma=4.5$ for the Median--MAD
rule. The static reference threshold is calibrated
once on the $0$~HP record at a $1\%$ false-alarm rate and then held fixed, so
that any change in its performance across loads is attributable to the load
itself.
